\section{Where the programme stands}

The programme has converged on a sharper diagnosis.  The principal missing
object is not another isolated local estimate.  It is a controlled middle
dynamics connecting ultraviolet gauge geometry to infrared spectral
structure.

\begin{principle}[Revised programme thesis]
The constructive Yang--Mills problem, in the present architecture, is to build
a gauge-covariant, positivity-preserving, topologically complete
renormalization dynamics whose normalized defect transfer is uniformly
contractive in physical units, and whose contraction is complete enough to
reconstruct a nonzero Hamiltonian spectral gap.
\end{principle}

This is stronger than saying that infrared rigidity must merely be
transported across representations.  Representation changes are not passive:
they forget information, alter the meaning of locality, and create lifting,
normalization, or exceptional-set defects.

The programme presently resembles a bridge with a developed ultraviolet
entrance, a conditional infrared endpoint, and a missing central span.

\paragraph{Ultraviolet and finite-block side.}
The package contains explicit \(\SU\) heat-kernel enclosures, quotient--fiber
certificates, nonlinear local charts, an exact Fr\'echet-dispersion identity,
dependent Gaussian calculations, boundary-frustration counterexamples,
state-weighted transfer lemmas, and partial center and integer topology
certificates.  These are restricted or finite statements.

\paragraph{Infrared and spectral side.}
The package contains an abstract representation calculus and conditional
theorems showing that a common temporal decay rate on a dense recovered
observable core excludes spectrum below that rate.  Countermodels show that
observable completeness, positivity, rate uniformity, and the order of
limits cannot be suppressed.

\paragraph{First missing physical implication.}
For every positive-mass retained history \(h\), one needs a normalized
comparison and chart-complement estimate of the form
\begin{align}
 \mu_h(A\cap G_h)&\le e^{\ell_h}\nu_h(A\cap G_h),\\
 \mu_h(G_h^c)&\le q_{\mathrm{exc},h},
\end{align}
such that
\begin{equation}
q_{\mathrm{low},h}+e^{\ell_h}q_{\mathrm{gauss},h}
+q_{\mathrm{exc},h}\le q_{\mathrm{cont}}
\end{equation}
uniformly in volume, cutoff, and admissible boundary state.  The package
closes this budget only in a centered reference model, not in the full
interacting measure.

\paragraph{Claim ceiling.}
The strongest Yang--Mills-facing status remains C2: rigorous restricted-model
theorems.  The abstract transfer-to-spectrum statements are conditional.
There is no physical one-scale basin theorem, repeated RG trajectory,
continuum construction, Osterwalder--Schrader reconstruction, or Hamiltonian
mass gap.

\subsection{The material orientation correction}

The orientation-correct Paper~05 result supersedes the exploratory Stage~44
zero-spill claim.  The centered covariance spectrum is
\begin{equation}
 \left\{0,\frac{239}{630},\frac{239}{630},
 \frac{33}{70},\frac{33}{70},\frac{33}{70}\right\},
\end{equation}
whereas the minimum external-curvature lift has strictly positive spectrum
\begin{equation}
 \left\{\frac{37}{33},\frac{37}{33},\frac{37}{33},
 \frac{391}{239},\frac{391}{239}\right\}.
\end{equation}
Every nonzero centered one-star displacement therefore spills curvature when
lifted to link variables.  The earlier zero mode resulted from centering raw
signed plaquette rows before putting all six staples in a common rooted-link
orientation.  This correction is evidence that typed representation maps are
part of the mathematics, not philosophical decoration.
